%% Beginning of file 'sample631.tex'
%%
%% Modified 2022 May  
%%
%% This is a sample manuscript marked up using the
%% AASTeX v6.31 LaTeX 2e macros.
%%
%% AASTeX is now based on Alexey Vikhlinin's emulateapj.cls 
%% (Copyright 2000-2015).  See the classfile for details.

%% AASTeX requires revtex4-1.cls and other external packages such as
%% latexsym, graphicx, amssymb, longtable, and epsf.  Note that as of 
%% Oct 2020, APS now uses revtex4.2e for its journals but remember that 
%% AASTeX v6+ still uses v4.1. All of these external packages should 
%% already be present in the modern TeX distributions but not always.
%% For example, revtex4.1 seems to be missing in the linux version of
%% TexLive 2020. One should be able to get all packages from www.ctan.org.
%% In particular, revtex v4.1 can be found at 
%% https://www.ctan.org/pkg/revtex4-1.

%% The first piece of markup in an AASTeX v6.x document is the \documentclass
%% command. LaTeX will ignore any data that comes before this command. The 
%% documentclass can take an optional argument to modify the output style.
%% The command below calls the preprint style which will produce a tightly 
%% typeset, one-column, single-spaced document.  It is the default and thus
%% does not need to be explicitly stated.
%%
%% using aastex version 6.3
\documentclass[linenumbers, twocolumn]{aastex631}

%% The default is a single spaced, 10 point font, single spaced article.
%% There are 5 other style options available via an optional argument. They
%% can be invoked like this:
%%
%% \documentclass[arguments]{aastex631}
%% 
%% where the layout options are:
%%
%%  twocolumn   : two text columns, 10 point font, single spaced article.
%%                This is the most compact and represent the final published
%%                derived PDF copy of the accepted manuscript from the publisher
%%  manuscript  : one text column, 12 point font, double spaced article.
%%  preprint    : one text column, 12 point font, single spaced article.  
%%  preprint2   : two text columns, 12 point font, single spaced article.
%%  modern      : a stylish, single text column, 12 point font, article with
%% 		  wider left and right margins. This uses the Daniel
%% 		  Foreman-Mackey and David Hogg design.
%%  RNAAS       : Supresses an abstract. Originally for RNAAS manuscripts 
%%                but now that abstracts are required this is obsolete for
%%                AAS Journals. Authors might need it for other reasons. DO NOT
%%                use \begin{abstract} and \end{abstract} with this style.
%%
%% Note that you can submit to the AAS Journals in any of these 6 styles.
%%
%% There are other optional arguments one can invoke to allow other stylistic
%% actions. The available options are:
%%
%%   astrosymb    : Loads Astrosymb font and define \astrocommands. 
%%   tighten      : Makes baselineskip slightly smaller, only works with 
%%                  the twocolumn substyle.
%%   times        : uses times font instead of the default
%%   linenumbers  : turn on lineno package.
%%   trackchanges : required to see the revision mark up and print its output
%%   longauthor   : Do not use the more compressed footnote style (default) for 
%%                  the author/collaboration/affiliations. Instead print all
%%                  affiliation information after each name. Creates a much 
%%                  longer author list but may be desirable for short 
%%                  author papers.
%% twocolappendix : make 2 column appendix.
%%   anonymous    : Do not show the authors, affiliations and acknowledgments 
%%                  for dual anonymous review.
%%
%% these can be used in any combination, e.g.
%%
%% \documentclass[twocolumn,linenumbers,trackchanges]{aastex631}
%%
%% AASTeX v6.* now includes \hyperref support. While we have built in specific
%% defaults into the classfile you can manually override them with the
%% \hypersetup command. For example,
%%
%% \hypersetup{linkcolor=red,citecolor=green,filecolor=cyan,urlcolor=magenta}
%%
%% will change the color of the internal links to red, the links to the
%% bibliography to green, the file links to cyan, and the external links to
%% magenta. Additional information on \hyperref options can be found here:
%% https://www.tug.org/applications/hyperref/manual.html#x1-40003
%%
%% Note that in v6.3 "bookmarks" has been changed to "true" in hyperref
%% to improve the accessibility of the compiled pdf file.
%%
%% If you want to create your own macros, you can do so
%% using \newcommand. Your macros should appear before
%% the \begin{document} command.
%%
\newcommand{\vdag}{(v)^\dagger}
\newcommand\aastex{AAS\TeX}
\newcommand\latex{La\TeX}

%% Reintroduced the \received and \accepted commands from AASTeX v5.2
%\received{March 1, 2021}
%\revised{April 1, 2021}
%\accepted{\today}

%% Command to document which AAS Journal the manuscript was submitted to.
%% Adds "Submitted to " the argument.
%\submitjournal{PSJ}

%% For manuscript that include authors in collaborations, AASTeX v6.31
%% builds on the \collaboration command to allow greater freedom to 
%% keep the traditional author+affiliation information but only show
%% subsets. The \collaboration command now must appear AFTER the group
%% of authors in the collaboration and it takes TWO arguments. The last
%% is still the collaboration identifier. The text given in this
%% argument is what will be shown in the manuscript. The first argument
%% is the number of author above the \collaboration command to show with
%% the collaboration text. If there are authors that are not part of any
%% collaboration the \nocollaboration command is used. This command takes
%% one argument which is also the number of authors above to show. A
%% dashed line is shown to indicate no collaboration. This example manuscript
%% shows how these commands work to display specific set of authors 
%% on the front page.
%%
%% For manuscript without any need to use \collaboration the 
%% \AuthorCollaborationLimit command from v6.2 can still be used to 
%% show a subset of authors.
%
%\AuthorCollaborationLimit=2
%
%% will only show Schwarz & Muench on the front page of the manuscript
%% (assuming the \collaboration and \nocollaboration commands are
%% commented out).
%%
%% Note that all of the author will be shown in the published article.
%% This feature is meant to be used prior to acceptance to make the
%% front end of a long author article more manageable. Please do not use
%% this functionality for manuscripts with less than 20 authors. Conversely,
%% please do use this when the number of authors exceeds 40.
%%
%% Use \allauthors at the manuscript end to show the full author list.
%% This command should only be used with \AuthorCollaborationLimit is used.

%% The following command can be used to set the latex table counters.  It
%% is needed in this document because it uses a mix of latex tabular and
%% AASTeX deluxetables.  In general it should not be needed.
%\setcounter{table}{1}

%%%%%%%%%%%%%%%%%%%%%%%%%%%%%%%%%%%%%%%%%%%%%%%%%%%%%%%%%%%%%%%%%%%%%%%%%%%%%%%%
%%
%% The following section outlines numerous optional output that
%% can be displayed in the front matter or as running meta-data.
%%
%% If you wish, you may supply running head information, although
%% this information may be modified by the editorial offices.
%\shorttitle{AASTeX v6.3.1 Sample article}
%\shortauthors{Schwarz et al.}
%%
%% You can add a light gray and diagonal water-mark to the first page 
%% with this command:
%% \watermark{text}
%% where "text", e.g. DRAFT, is the text to appear.  If the text is 
%% long you can control the water-mark size with:
%% \setwatermarkfontsize{dimension}
%% where dimension is any recognized LaTeX dimension, e.g. pt, in, etc.
%%
%%%%%%%%%%%%%%%%%%%%%%%%%%%%%%%%%%%%%%%%%%%%%%%%%%%%%%%%%%%%%%%%%%%%%%%%%%%%%%%%
%\graphicspath{{./}{figures/}}
%% This is the end of the preamble.  Indicate the beginning of the
%% manuscript itself with \begin{document}.

\begin{document}

\title{Color-Magnitude and Hertzsprung-Russell Diagram Morphology of Simulated Star Cluster Formation}

%% LaTeX will automatically break titles if they run longer than
%% one line. However, you may use \\ to force a line break if
%% you desire. In v6.31 you can include a footnote in the title.

%% A significant change from earlier AASTEX versions is in the structure for 
%% calling author and affiliations. The change was necessary to implement 
%% auto-indexing of affiliations which prior was a manual process that could 
%% easily be tedious in large author manuscripts.
%%
%% The \author command is the same as before except it now takes an optional
%% argument which is the 16 digit ORCID. The syntax is:
%% \author[xxxx-xxxx-xxxx-xxxx]{Author Name}
%%
%% This will hyperlink the author name to the author's ORCID page. Note that
%% during compilation, LaTeX will do some limited checking of the format of
%% the ID to make sure it is valid. If the "orcid-ID.png" image file is 
%% present or in the LaTeX pathway, the OrcID icon will appear next to
%% the authors name.
%%
%% Use \affiliation for affiliation information. The old \affil is now aliased
%% to \affiliation. AASTeX v6.31 will automatically index these in the header.
%% When a duplicate is found its index will be the same as its previous entry.
%%
%% Note that \altaffilmark and \altaffiltext have been removed and thus 
%% can not be used to document secondary affiliations. If they are used latex
%% will issue a specific error message and quit. Please use multiple 
%% \affiliation calls for to document more than one affiliation.
%%
%% The new \altaffiliation can be used to indicate some secondary information
%% such as fellowships. This command produces a non-numeric footnote that is
%% set away from the numeric \affiliation footnotes.  NOTE that if an
%% \altaffiliation command is used it must come BEFORE the \affiliation call,
%% right after the \author command, in order to place the footnotes in
%% the proper location.
%%
%% Use \email to set provide email addresses. Each \email will appear on its
%% own line so you can put multiple email address in one \email call. A new
%% \correspondingauthor command is available in V6.31 to identify the
%% corresponding author of the manuscript. It is the author's responsibility
%% to make sure this name is also in the author list.
%%
%% While authors can be grouped inside the same \author and \affiliation
%% commands it is better to have a single author for each. This allows for
%% one to exploit all the new benefits and should make book-keeping easier.
%%
%% If done correctly the peer review system will be able to
%% automatically put the author and affiliation information from the manuscript
%% and save the corresponding author the trouble of entering it by hand.

%\correspondingauthor{August Muench}
%\email{greg.schwarz@aas.org, gus.muench@aas.org}

\author[0009-0004-8200-9910]{Winston Zhang}
\affiliation{University of Virginia, Department of Astronomy \\
Charlottesville, VA 22904, USA}

\author[0000-0002-4725-7589]{Aayush Gautam}
\affiliation{Tribhuvan University, Department of Physics \\
Bharatpur 44207, Nepal}
\affiliation{University of Virginia, Department of Astronomy \\
Charlottesville, VA 22904, USA}

\author[0000-0002-5851-2602]{Juan Farias}
\affiliation{Chalmers University of Technology, Department of Space, Earth \& Environment \\
Gothenburg, Sweden}
\affiliation{University of Texas at Austin, Department of Astronomy \\
Austin, TX 78712, USA}

%\author{Aabha Brizawasi}
%\affiliation{Indian Institute of Science Education and Research \\
% Mohali, India}

\author{Morten Andersen}
\affiliation{European Southern Observatory \\
München, Germany}

\author{Jonathan Tan}
\affiliation{University of Virginia, Department of Astronomy \\
Charlottesville, VA 22904, USA}
\affiliation{Chalmers University of Technology, Department of Space, Earth \& Environment \\
Gothenburg, Sweden}

%% Note that the \and command from previous versions of AASTeX is now
%% depreciated in this version as it is no longer necessary. AASTeX 
%% automatically takes care of all commas and "and"s between authors names.

%% AASTeX 6.31 has the new \collaboration and \nocollaboration commands to
%% provide the collaboration status of a group of authors. These commands 
%% can be used either before or after the list of corresponding authors. The
%% argument for \collaboration is the collaboration identifier. Authors are
%% encouraged to surround collaboration identifiers with ()s. The 
%% \nocollaboration command takes no argument and exists to indicate that
%% the nearby authors are not part of surrounding collaborations.

%% Mark off the abstract in the ``abstract'' environment. 
\begin{abstract}

placeholder abstract

\end{abstract}

%% Keywords should appear after the \end{abstract} command. 
%% The AAS Journals now uses Unified Astronomy Thesaurus concepts:
%% https://astrothesaurus.org
%% You will be asked to selected these concepts during the submission process
%% but this old "keyword" functionality is maintained in case authors want
%% to include these concepts in their preprints.
\keywords{}

%% From the front matter, we move on to the body of the paper.
%% Sections are demarcated by \section and \subsection, respectively.
%% Observe the use of the LaTeX \label
%% command after the \subsection to give a symbolic KEY to the
%% subsection for cross-referencing in a \ref command.
%% You can use LaTeX's \ref and \label commands to keep track of
%% cross-references to sections, equations, tables, and figures.
%% That way, if you change the order of any elements, LaTeX will
%% automatically renumber them.
%%
%% We recommend that authors also use the natbib \citep
%% and \citet commands to identify citations.  The citations are
%% tied to the reference list via symbolic KEYs. The KEY corresponds
%% to the KEY in the \bibitem in the reference list below. 

\section{Introduction} \label{sec:intro}
placeholder intro





\section{Related Work}\label{sec:related}
The formation and early evolution of star clusters have been explored extensively through numerical simulations that couple gravitational dynamics with prescriptions for star formation. In particular, \citet{Farias2019} present Nbody6 simulations of protoclusters forming within collapsing molecular clouds, demonstrating that the cloud surface density and star formation efficiency per freefall time ($\epsilon_{\mathrm{ff}}$) regulate both the duration of star formation and the resulting stellar age spread. These simulations provide a physically motivated link between cloud-scale properties and observable stellar populations.

Synthetic diagram modeling has previously been used to connect such simulations to observable parameter spaces. \citet{Masegian2021} generated synthetic Hertzsprung–Russell (H–R) diagrams by interpolating stellar masses and ages from cluster simulations onto theoretical isochrones, quantifying the relationship between age spread and luminosity dispersion. That work established that luminosity spread in H–R space is positively correlated with the underlying stellar age spread induced by extended star formation histories.

Extending this approach into observable photometric space, \citet{Zhang2024} implemented a forward-modeling framework that combined Nbody6 simulation outputs with synthetic photometry from SPISEA \citep{Hosek2020} to generate color–magnitude diagrams (CMDs) of forming protoclusters. That study demonstrated that magnitude spread in CMD space similarly traces stellar age spread and therefore encodes information about the parent cloud’s collapse timescale and density.

Together, these prior works establish a forward relationship between cloud physics, stellar age spread, and diagram morphology. The present study builds upon these foundations by further developing the synthetic CMD framework, examining its behavior across multiple photometric representations, and quantifying how CMD morphology reflects the underlying star formation history of a protocluster.

\section{Methods} \label{sec:methods}

We develop a forward-modeling framework that links protocluster formation simulations to observable color--magnitude diagrams (CMDs) and Hertzsprung--Russell (H--R) diagrams. The methodology combines numerical cluster simulations with stellar-evolution models and synthetic photometry in order to map simulated stellar masses and ages into theoretical and observable diagram spaces. We then quantify the morphology of the resulting stellar sequences and examine how that morphology depends on the star-formation history of the parent protocluster.

The forward-modeling pipeline combines protocluster simulations from \citet{Farias2019} with synthetic photometry generated using SPISEA \citep{Hosek2020}. Stellar masses and ages extracted from Nbody6 simulations are interpolated onto a grid of theoretical isochrones. The resulting stellar populations are analyzed both in H--R space, using effective temperature and luminosity, and in several photometric CMD projections.

\subsection{Protocluster Simulations}

We utilize Nbody6 simulations presented in \citet{Farias2019}, which model the gradual formation of embedded star clusters within collapsing molecular clouds. The present analysis is restricted to the simulation suite with an initial clump mass of $M_{\rm cl}=3000\,M_\odot$. These simulations adopt a Kroupa initial mass function and include primordial binary systems.

The primary cloud-scale parameters considered here are the cloud surface density $\Sigma_{\rm cloud}$ and the star formation efficiency per free-fall time $\epsilon_{\rm ff}$. The parameter $\epsilon_{\rm ff}$ regulates how rapidly stars are introduced relative to the free-fall time of the parent cloud. Lower values therefore produce more extended star-formation histories, while larger values approach increasingly rapid formation. The cloud surface density also changes the dynamical and formation timescales of the protocluster and consequently affects the stellar age distribution.

At a specified cluster time $t_{\rm cl}$, an Nbody6 snapshot is converted into a stellar population containing the masses and ages of the stars present at that epoch. These quantities form the input to the stellar-evolution and synthetic-photometry interpolation described below.

\subsection{Simulation Parameter Grid}
\label{sec:simulation-grid}

For the adopted initial clump mass of $M_{\rm cl}=3000\,M_\odot$, the simulation suite spans two cloud surface densities,
\begin{equation}
\Sigma_{\rm cloud}
=
0.1,\ 1.0\ {\rm g\,cm^{-2}},
\end{equation}
and five values of the star formation efficiency per free-fall time,
\begin{equation}
\epsilon_{\rm ff}
=
0.01,\ 0.03,\ 0.10,\ 0.30,\ 1.00.
\end{equation}

Each combination of $\Sigma_{\rm cloud}$ and $\epsilon_{\rm ff}$ is represented by 20 independent realizations, identified by seed numbers 00--19. The full analysis therefore contains 200 finite-formation simulations.

For a fixed cloud surface density, simulations with the same seed but different $\epsilon_{\rm ff}$ values contain the same underlying set of stellar masses sampled from the initial mass function. Changing $\epsilon_{\rm ff}$ for a fixed seed primarily changes the times at which these stars are introduced into the cluster. This enables controlled same-seed comparisons in which differences in stellar population morphology can be associated with the star-formation history rather than with a different random realization of the initial mass function.

For ensemble measurements, all 20 realizations of a given $(\Sigma_{\rm cloud},\epsilon_{\rm ff})$ combination are analyzed independently. We report the arithmetic mean across the contributing seeds, while the realization-to-realization scatter is characterized by the sample standard deviation. The resulting standard-deviation bands therefore represent stochastic variation between cluster realizations rather than uncertainty in the ensemble mean.

\begin{deluxetable}{ccc}
\tablecaption{$3000\,M_\odot$ Protocluster Simulation Grid
\label{tab:simulation-grid}}
\tablehead{
\colhead{$\Sigma_{\rm cloud}$ (${\rm g\,cm^{-2}}$)} &
\colhead{$\epsilon_{\rm ff}$} &
\colhead{Number of realizations}
}
\startdata
0.1 & 0.01 & 20 \\
0.1 & 0.03 & 20 \\
0.1 & 0.10 & 20 \\
0.1 & 0.30 & 20 \\
0.1 & 1.00 & 20 \\
1.0 & 0.01 & 20 \\
1.0 & 0.03 & 20 \\
1.0 & 0.10 & 20 \\
1.0 & 0.30 & 20 \\
1.0 & 1.00 & 20 \\
\enddata
\end{deluxetable}

All results presented in this work are based on the $M_{\rm cl}=3000\,M_\odot$ simulation suite. Dependence on initial clump mass is reserved for future analysis.

\subsection{Isochrone Grid Construction}

Synthetic stellar populations are generated using SPISEA \citep{Hosek2020}. The analysis adopts the merged Baraffe--Pisa--Ekstr\"om--PARSEC evolutionary model grid used by the forward-modeling pipeline. Synthetic photometry is calculated with the BT-Settl 2015 stellar-atmosphere models.

The isochrone grid spans ages from 1 to 20 Myr in intervals of 0.5 Myr. We assume solar metallicity, zero extinction ($A_{K_s}=0$), and a distance of 410 pc. The same stellar-evolution and photometric assumptions are held fixed throughout the simulation ensemble so that changes in diagram morphology arise from changes in the underlying simulated stellar population rather than from changes in the synthetic-observation model.

Photometry is generated in the JWST F070W, F182M, and F200W filters and the HST F555W and F814W filters. These bands permit construction of several CMD projections with substantially different isochrone geometries.

\subsection{Photometric Interpolation}

For each simulated star with mass $m$ and age $a$, stellar parameters and photometric magnitudes are obtained by interpolation within the isochrone grid. For ages lying between two grid points, the procedure is:

\begin{enumerate}
    \item identify the two isochrones whose ages bracket $a$;
    \item interpolate in stellar mass along each bracketing isochrone;
    \item interpolate between the two mass-interpolated points in age.
\end{enumerate}

The resulting stellar properties include effective temperature, luminosity, surface gravity, and magnitudes in all selected filters. Stars younger than the minimum age represented by the stellar-evolution grid are evaluated at the youngest available grid age, while stars outside the supported upper age range are excluded.

The interpolation procedure is illustrated schematically in Figure~\ref{fig:iso-interp-example}.

\begin{figure}[ht]
    \centering
    \includegraphics[width=1\linewidth]{science-images/fig_interpolation_schematic_f162m_f182m.png}
    \caption{Example bilinear interpolation using neighboring stellar isochrones to obtain the synthetic properties of a star whose mass and age do not coincide exactly with the discrete isochrone grid.}
    \label{fig:iso-interp-example}
\end{figure}

\subsection{Treatment of Binary Systems}

The Nbody6 simulations include primordial binary systems. Binary stars may be represented in two limiting observational cases.

\paragraph{Resolved binaries:}
The primary and companion are treated as independent stellar sources and are interpolated separately.

\paragraph{Unresolved binaries:}
The fluxes of the two components are summed before conversion back into a magnitude,
\begin{equation}
F_{\rm binary}=F_1+F_2,
\end{equation}
and
\begin{equation}
m_{\rm binary}
=
-2.5\log_{10}
\left(
\frac{F_{\rm binary}}{F_0}
\right),
\end{equation}
where $F_0$ is the appropriate photometric zero-point flux.

Unresolved binaries therefore appear brighter than either component individually and can introduce a secondary sequence or additional vertical width in CMD space. Figure~\ref{fig:cmd-binaries} illustrates the difference between these limiting treatments.

\begin{figure*}
    \centering
    \includegraphics[width=1\linewidth]{science-images/fig_cmd_binaries_f162m_f182m.png}
    \caption{Example synthetic CMDs demonstrating the effect of treating simulated binary systems as resolved individual sources or unresolved combined sources.}
    \label{fig:cmd-binaries}
\end{figure*}

\subsection{Synthetic Diagram Construction}

For every interpolated stellar population, we construct both theoretical H--R diagrams and synthetic CMDs. The H--R representation uses $T_{\rm eff}$ versus $\log_{10}(L/L_\odot)$,
while photometric colors are formed from differences between synthetic magnitudes.

The complete ensemble analysis considers the following three CMD spaces: (F070W-F200W) versus F200W, (F182M-F200W) versus F200W, and (F555W-F814W. versus F814W.

Together with the H--R diagram, these projections allow us to determine whether trends present in physical stellar-parameter space remain visible after transformation into observable photometric space.

\subsection{Definition of the Measurement Region}
\label{sec:measurement-region}

The morphology statistic is measured over a restricted portion of the low-mass pre--main-sequence population. A star is eligible for the diagram-spread analysis only when 3000 K $\leq T_{ eff} \leq$ 3600 K. The measurement geometry is fixed for all cluster ages and all simulations.

In H--R space, the interval from 3000 to 3600 K is divided directly into six 100 K bins,
\begin{equation}
3000,\ 3100,\ \ldots,\ 3600\ {\rm K}.
\end{equation}

The corresponding CMD binning is defined differently because color is a nonlinear function of effective temperature. For each CMD projection, the 1 Myr reference isochrone is used to determine the color corresponding to $T_{\rm eff}=3000$ K and the color corresponding to $T_{\rm eff}=3600$ K. The fixed color interval between these two values is then divided into six equal-width bins. These color boundaries are held fixed for every subsequent simulation snapshot and cluster age. Individual stars are assigned to the bins according to their actual interpolated color.

This procedure provides a common low-mass stellar interval across the different diagram spaces while avoiding age-dependent movement of the bin boundaries. Figure~\ref{fig:analysis-geometry} illustrates the isochrone geometry, representative stellar population, and fixed measurement boundaries used in each diagram.

\begin{figure*}
    \centering
    \includegraphics[width=1\linewidth]{science-images/fig_analysis_space_isochrone_geometry.png}
    \caption{Isochrone geometry, fixed measurement bins, and a representative simulated stellar population in the four diagram spaces used for the ensemble analysis. Isochrones are colored by age. In H--R space the boundaries correspond to six 100 K intervals between 3000 and 3600 K. In each CMD, the boundaries are six equal-width color bins defined between the colors corresponding to 3000 and 3600 K on the 1 Myr reference isochrone.}
    \label{fig:analysis-geometry}
\end{figure*}

\subsection{Diagram Morphology Quantification}
\label{sec:spread-metric}

We quantify the vertical extent of the stellar sequence using the separation between the lower and upper quartiles of the distribution in each horizontal bin. The method follows the luminosity-spread approach of \citet{Masegian2021}, adapted here for both H--R and CMD space.

For a valid horizontal bin $j$, the 25th and 75th percentiles of the vertical coordinate are first calculated. We then compute the mean vertical coordinate of the stars in the lower quartile and the mean vertical coordinate of the stars in the upper quartile.

For a CMD, where the vertical coordinate is magnitude, the bin spread is
\begin{equation}
S_{m,j}
=
\left\langle m\right\rangle_{{\rm faint},j}
-
\left\langle m\right\rangle_{{\rm bright},j},
\end{equation}
such that a broader stellar sequence gives a larger positive value.

For an H--R diagram, the analogous luminosity spread is
\begin{equation}
S_{L,j}
=
\left\langle \log(L/L_\odot)\right\rangle_{{\rm upper},j}
-
\left\langle \log(L/L_\odot)\right\rangle_{{\rm lower},j}.
\end{equation}

A bin must contain at least eight eligible stars to contribute to the measurement. The overall spread is the arithmetic mean over all valid bins,
\begin{equation}
S
=
\frac{1}{N_{\rm valid}}
\sum_{j=1}^{N_{\rm valid}}S_j.
\end{equation}

Using the means of the outer quartiles suppresses sensitivity to individual outliers while retaining sensitivity to broadening of the stellar sequence. The fixed six-bin construction also ensures that the same portion of diagram space is sampled consistently across simulations and cluster ages.

\subsection{Intrinsic Stellar Age Spread}

The intrinsic stellar age spread of each finite-formation simulation snapshot is quantified independently of the diagram morphology as
\begin{equation}
\Delta t_{90}
=
t_{95}-t_5,
\end{equation}
where $t_{95}$ and $t_5$ are the 95th and 5th percentiles of the stellar-age distribution.

The age percentiles are calculated from the full successfully interpolated stellar population in the snapshot, rather than only from the 3000--3600 K subset used for the diagram-spread measurement. This definition suppresses sensitivity to a small number of extreme-age stars while characterizing the duration represented by the bulk of the stellar population.

\subsection{Ensemble Analysis}

The spread measurement is evaluated at cluster times from 1 to 20 Myr in 0.5 Myr intervals for every available realization, yielding a time-dependent measurement for each $(\Sigma_{\rm cloud},\epsilon_{\rm ff},{\rm seed})$ combination and diagram space.

For each $(\Sigma_{\rm cloud},\epsilon_{\rm ff})$ family, the measurements from the 20 seeds are combined to determine the ensemble mean and sample standard deviation at each cluster time. We compare these ensemble measurements in three complementary ways:

\begin{enumerate}
    \item the evolution of the intrinsic stellar age spread $\Delta t_{90}$ with cluster time;
    \item the evolution of diagram spread with cluster time;
    \item the direct relation between diagram spread and intrinsic stellar age spread.
\end{enumerate}

This ensemble treatment separates systematic trends with the cloud parameters from stochastic differences between individual realizations and allows the dependence of diagram morphology on $\epsilon_{\rm ff}$ and $\Sigma_{\rm cloud}$ to be evaluated across the complete simulation grid.


\section{Results \& Analysis} \label{sec:results}

The forward-modeling framework successfully maps the simulated stellar populations into both H--R and color--magnitude space. The resulting diagrams evolve systematically with the stellar age distribution and retain measurable signatures of the cloud parameters that regulate cluster formation.

\subsection{Synthetic Cluster Morphology}

Figure~\ref{fig:representative-cluster} shows a representative simulated cluster in the $(F182M-F200W)$ versus F200W CMD together with the synthetic isochrone grid. Individual simulated stars populate a broad pre--main-sequence locus bounded by isochrones of different ages. The population does not lie on a single-age sequence because star formation proceeds over a finite interval in the simulations.

\begin{figure}
    \centering
    \includegraphics[width=1\linewidth]{science-images/fig_representative_cluster_f182m_f200w.png}
    \caption{Representative synthetic cluster in $(F182M-F200W)$ versus F200W space, overlaid on the 1--20 Myr isochrone grid. Isochrones are colored by age using the same color scale adopted throughout the diagram-geometry figures.}
    \label{fig:representative-cluster}
\end{figure}

The morphology of an individual cluster also changes with cluster time. Figure~\ref{fig:cmd-age-evolution} follows one realization at six epochs in the $(F182M-F200W)$ versus F200W projection. The same isochrone grid is shown in each panel, allowing the evolution of the population to be separated visually from changes in the underlying reference tracks.

At early times, the population contains a substantial fraction of very young, luminous pre--main-sequence stars. As the cluster ages, the entire sequence evolves toward fainter magnitudes while the stellar age distribution also changes as star formation proceeds and ultimately subsides. The resulting CMD width therefore reflects both stellar evolution and the duration of star formation.

\begin{figure*}
    \centering
    \includegraphics[width=1\linewidth]{science-images/fig_cluster_time_evolution_f182m_f200w.png}
    \caption{Evolution of one simulated cluster realization in $(F182M-F200W)$ versus F200W space. Each panel shows the same 1--20 Myr isochrone grid and the interpolated stellar population at a different cluster time.}
    \label{fig:cmd-age-evolution}
\end{figure*}

\subsection{Dependence on Diagram Space}

The same stellar population can have substantially different apparent morphology depending on the coordinate system in which it is viewed. Figure~\ref{fig:analysis-geometry} compares the H--R diagram to the three CMD projections included in the ensemble analysis.

In H--R space, the low-mass pre--main-sequence tracks remain relatively well ordered in effective temperature and luminosity. The transformation into photometric space alters both the slope and curvature of these tracks. In particular, the same $3000$--$3600$ K stellar interval can span very different ranges in color depending on the adopted filter pair.

This motivates the use of a separately defined fixed color interval for each CMD rather than treating equal color intervals in different filters as physically equivalent. The $(F182M-F200W)$ versus F200W projection provides a comparatively well-behaved mapping of the selected temperature range and remains a useful principal photometric representation, while the additional JWST and HST projections provide tests of how robust the inferred trends are to the choice of observable space.

\subsection{Evolution of the Intrinsic Stellar Age Spread}

Figure~\ref{fig:age-spread-vs-time} shows the ensemble-averaged intrinsic age spread as a function of cluster time for both cloud surface densities and all five values of $\epsilon_{\rm ff}$. Each curve represents the mean across the available seeds, and the shaded region represents the sample standard deviation among realizations.

A clear ordering with $\epsilon_{\rm ff}$ emerges. Simulations with lower $\epsilon_{\rm ff}$ develop larger stellar age spreads because star formation proceeds over a longer interval. Conversely, the highest-efficiency simulations form their stellar populations more rapidly and consequently maintain substantially narrower age distributions.

The cloud surface density introduces a second systematic dependence. At fixed $\epsilon_{\rm ff}$, the $\Sigma_{\rm cloud}=0.1\,{\rm g\,cm^{-2}}$ simulations generally develop larger age spreads than the $\Sigma_{\rm cloud}=1.0\,{\rm g\,cm^{-2}}$ simulations. The lower-density clouds therefore produce more temporally extended stellar populations in the simulation suite considered here.

At early cluster times, $\Delta t_{90}$ grows as newly formed stars continue to extend the age distribution. Once star formation has largely ceased, the difference between the oldest and youngest major components of the population approaches a characteristic value and the age-spread curves flatten.

\begin{figure*}
    \centering
    \includegraphics[width=1\linewidth]{science-images/fig_ensemble_age_spread_vs_time.png}
    \caption{Ensemble intrinsic stellar age spread $\Delta t_{90}$ as a function of cluster time. Separate panels show $\Sigma_{\rm cloud}=0.1$ and $1.0\,{\rm g\,cm^{-2}}$. Curves show the mean across the 20 simulation realizations and shaded regions show the sample standard deviation among seeds. Lower $\epsilon_{\rm ff}$ produces systematically larger age spreads, while the lower cloud-surface-density simulations also exhibit larger age spreads at fixed $\epsilon_{\rm ff}$.}
    \label{fig:age-spread-vs-time}
\end{figure*}

\subsection{Diagram Spread as a Function of Cluster Time}

The same hierarchy appears in the morphology of the synthetic stellar diagrams. Figure~\ref{fig:spread-vs-time} shows the ensemble quartile-tail spread as a function of cluster time in H--R space and in all three CMD projections.

Across the diagram spaces, lower values of $\epsilon_{\rm ff}$ generally produce larger measured spreads. This mirrors the intrinsic age-spread hierarchy and demonstrates that the extended star-formation histories of the low-efficiency models remain encoded in the morphology of the stellar sequence after transformation into photometric space.

Cloud surface density produces a corresponding systematic effect. For a fixed $\epsilon_{\rm ff}$, clusters formed in the lower-surface-density $\Sigma_{\rm cloud}=0.1\,{\rm g\,cm^{-2}}$ models generally exhibit larger diagram spreads than their $\Sigma_{\rm cloud}=1.0\,{\rm g\,cm^{-2}}$ counterparts. Thus the ordering identified directly in the stellar-age distributions is also recovered from the H--R and CMD morphology.

The absolute value and detailed temporal behavior of the spread depend on diagram space. This is expected because the stellar-evolution tracks are transformed nonlinearly into different photometric coordinates. The agreement in the relative ordering of the simulation families is therefore more physically informative than equality of the numerical spread measured in different diagrams.

Some temporal structure also remains after averaging over 20 realizations. Diagram spread is not simply a monotonic function of cluster time because the morphology is affected simultaneously by the development of the age distribution and by the evolution of the pre--main-sequence isochrones themselves. Nevertheless, the ensemble analysis reveals a persistent separation between simulation families with different formation histories.

\begin{figure*}
    \centering
    \includegraphics[width=1\linewidth]{science-images/fig_ensemble_spread_vs_time_all_spaces.png}
    \caption{Ensemble quartile-tail diagram spread as a function of cluster time in H--R space and the three analyzed CMD projections. Columns correspond to the two cloud surface densities. Finite-formation curves show the mean across simulation seeds, with shaded regions indicating the sample standard deviation. Lower $\epsilon_{\rm ff}$ and lower $\Sigma_{\rm cloud}$ produce systematically broader stellar sequences.}
    \label{fig:spread-vs-time}
\end{figure*}

\subsection{Diagram Spread as a Tracer of Intrinsic Age Spread}

The connection between the observable morphology and the underlying stellar-age distribution is shown more directly in Figure~\ref{fig:spread-vs-age-spread}. Rather than plotting both quantities against cluster time, the ensemble mean diagram spread is plotted directly against the corresponding ensemble mean $\Delta t_{90}$.

In H--R space, increasing intrinsic age spread is accompanied by increasing luminosity spread over much of the simulation sequence. This reproduces the qualitative behavior identified by \citet{Masegian2021}: a population containing stars formed over a longer interval occupies a broader vertical region of the pre--main-sequence H--R diagram.

The same relationship remains visible after transformation into CMD space. In particular, the $(F182M-F200W)$ versus F200W projection retains a clear connection between stellar age spread and photometric sequence width. The F070W--F200W and HST F555W--F814W projections show that the detailed shape and normalization of the relation depend on filter choice, but the underlying tendency for larger intrinsic age spreads to produce broader diagrams remains present.

The simulation families occupy different regions of this relation because both $\epsilon_{\rm ff}$ and $\Sigma_{\rm cloud}$ regulate the duration of star formation. Low-$\epsilon_{\rm ff}$ models preferentially occupy the high-$\Delta t_{90}$, high-spread portion of the diagrams, whereas high-$\epsilon_{\rm ff}$ models are concentrated toward smaller age spreads and narrower sequences. Likewise, the lower-density cloud models are displaced toward larger age and diagram spreads relative to the higher-density models.

\begin{figure*}
    \centering
    \includegraphics[width=1\linewidth]{science-images/fig_ensemble_spread_vs_age_spread_all_spaces.png}
    \caption{Ensemble diagram spread as a function of intrinsic stellar age spread $\Delta t_{90}$ for the H--R diagram and three CMD projections. Each sequence traces the evolution of a simulation family through cluster time. Larger intrinsic age spreads generally correspond to broader stellar sequences, with low-$\epsilon_{\rm ff}$ and low-$\Sigma_{\rm cloud}$ models preferentially occupying the high-spread region. Error bars indicate the realization-to-realization sample standard deviation.}
    \label{fig:spread-vs-age-spread}
\end{figure*}

\subsection{Dependence on Star Formation Efficiency and Cloud Surface Density}

The ensemble results demonstrate that the diagram spread is sensitive to both principal cloud parameters explored in the simulation suite.

The strongest systematic ordering is associated with $\epsilon_{\rm ff}$. Lower star formation efficiencies per free-fall time correspond to more prolonged formation histories. These simulations therefore contain stars spanning a larger range of formation times and develop larger values of $\Delta t_{90}$. The increased age diversity broadens the pre--main-sequence population vertically in H--R space and produces a corresponding increase in magnitude spread in the synthetic CMDs.

Cloud surface density modifies the same relationship. At otherwise fixed $\epsilon_{\rm ff}$, the lower-density $\Sigma_{\rm cloud}=0.1\,{\rm g\,cm^{-2}}$ simulations exhibit larger intrinsic age spreads and larger diagram spreads than the $\Sigma_{\rm cloud}=1.0\,{\rm g\,cm^{-2}}$ simulations. The synthetic diagrams therefore retain information not only about the duration of star formation as controlled by $\epsilon_{\rm ff}$, but also about the physical conditions of the parent cloud.

These trends are recovered after averaging over the full 20-seed ensembles and are therefore not attributable to a particular stochastic sampling of the stellar initial mass function. Individual realizations exhibit significant scatter, particularly at some cluster ages and in some photometric projections, but the ordering of the ensemble means remains systematic.

\subsection{Implications for Synthetic CMD Diagnostics}

The ensemble analysis establishes a forward connection between cloud-scale star-formation parameters, intrinsic stellar age distributions, and observable diagram morphology. In the simulation suite considered here, we find that a lower freefall parameter $\epsilon_{\rm ff}$ (i.e., lower star-forming efficient) and lower cloud density $\Sigma_{cloud}$ both independently contribute to a larger age spread $\Delta_{90}$, and transitively a larger diagram spread.

The correspondence is clearest when interpreted statistically across the simulation ensemble rather than from an individual snapshot. Stochastic differences between realizations and changes in isochrone geometry prevent the diagram width from serving as a unique one-to-one estimator of age spread at every epoch. Nevertheless, the persistence of the parameter ordering across H--R and multiple CMD projections shows that the width of the pre--main-sequence population contains measurable information about the underlying formation history.

The forward-modeling framework therefore provides a basis for connecting observed CMD morphology to physically motivated cluster-formation models. A subsequent observational comparison must additionally account for extinction, photometric uncertainty, completeness, unresolved multiplicity, and field contamination, but the simulation analysis presented here establishes the intrinsic relationship that such observational effects would act upon.





\section{Conclusion} \label{sec:conclusion}
placeholder conclusion

%% IMPORTANT! The old "\acknowledgment" command has be depreciated. It was
%% not robust enough to handle our new dual anonymous review requirements and
%% thus been replaced with the acknowledgment environment. If you try to 
%% compile with \acknowledgment you will get an error print to the screen
%% and in the compiled pdf.
%% 
%% Also note that the akcnowlodgment environment does not support long amounts of text. If you have a lot of people and institutions to acknowledge, do not use this command. Instead, create a new \section{Acknowledgments}.
\begin{acknowledgments}

\end{acknowledgments}

%% To help institutions obtain information on the effectiveness of their 
%% telescopes the AAS Journals has created a group of keywords for telescope 
%% facilities.
%
%% Following the acknowledgments section, use the following syntax and the
%% \facility{} or \facilities{} macros to list the keywords of facilities used 
%% in the research for the paper.  Each keyword is check against the master 
%% list during copy editing.  Individual instruments can be provided in 
%% parentheses, after the keyword, but they are not verified.

\software{astropy \citep{2013A&A...558A..33A,2018AJ....156..123A}
          }

%% Appendix material should be preceded with a single \appendix command.
%% There should be a \section command for each appendix. Mark appendix
%% subsections with the same markup you use in the main body of the paper.

%% Each Appendix (indicated with \section) will be lettered A, B, C, etc.
%% The equation counter will reset when it encounters the \appendix
%% command and will number appendix equations (A1), (A2), etc. The
%% Figure and Table counter will not reset.

\appendix

\section{Appendix information}

Appendices can be broken into separate sections just like in the main text.
The only difference is that each appendix section is indexed by a letter
(A, B, C, etc.) instead of a number.  Likewise numbered equations have
the section letter appended.  Here is an equation as an example.
\begin{equation}
I = \frac{1}{1 + d_{1}^{P (1 + d_{2} )}}
\end{equation}
Appendix tables and figures should not be numbered like equations. Instead
they should continue the sequence from the main article body.

%% For this sample we use BibTeX plus aasjournals.bst to generate the
%% the bibliography. The sample631.bib file was populated from ADS. To
%% get the citations to show in the compiled file do the following:
%%
%% pdflatex sample631.tex
%% bibtext sample631
%% pdflatex sample631.tex
%% pdflatex sample631.tex

\bibliography{sources}{}
\bibliographystyle{aasjournal}

%% This command is needed to show the entire author+affiliation list when
%% the collaboration and author truncation commands are used.  It has to
%% go at the end of the manuscript.
%\allauthors

%% Include this line if you are using the \added, \replaced, \deleted
%% commands to see a summary list of all changes at the end of the article.
%\listofchanges

\end{document}

% End of file `sample631.tex'.
